% GENERATED FILE. Edit canonical lessons, then run npm run generate:books.
% canonical-source-hash: fnv1a64:cf57d308a6975f38
% canonical-lessons: IT-C516-violoncello, IT-C516-flauto, IT-C516-tromba, IT-C516-fisarmonica, IT-C516-strimpellare

\chapter{Strumenti musicali}
\label{ch:it-strumenti-musicali}
\hlchaptermodality{516}

\begin{chapteropening}
\textbf{By the end of this chapter:} \emph{I can talk about the instruments people play.}

Complete the last lesson of chapter 516: I can talk about the instruments people play.
\end{chapteropening}

\section[il violoncello]{il violoncello — cello}

\label{lesson:IT-C516-violoncello}



\begin{quote}
Before the new one: say the Italian for skier, then the Italian for chairlift.
\end{quote}

\subsection*{You'll want to know: il violoncello}
\textbf{il violoncello} — \textquotedblleft{}cello\textquotedblright{}.

\begin{quote}
Suono il violoncello da dieci anni.
\end{quote}

\begin{quote}
\emph{I have played the cello for ten years.}
\end{quote}

Often shortened to \emph{il cello}.

\begin{cousinweb}
\emph{violoncello} is built on \emph{violone} \textquotedblleft{}a large viol\textquotedblright{} with \emph{-cello}, a little ending (EtymDB): a smaller version of the big viol.
\end{cousinweb}

\begin{grammarlens}[title={What you've built}]
One more word for telling, explaining and describing at B1.
\end{grammarlens}

\subsection*{Your turn}
\begin{itemize}
\raggedright
  \item \emph{Say it:} \emph{il violoncello}
  \item \emph{Say it:} \emph{il violoncello}, once more
  \item \emph{Say it:} a sentence of your own with \emph{il violoncello}
  \item \emph{Recall:} say the Italian for skier, then the Italian for chairlift, then say \emph{il violoncello} again
\end{itemize}

\begin{tcolorbox}[breakable,colback=teal!4,colframe=teal!35!black,title={Before you move on}]
What does il violoncello mean? (\textquotedblleft{}Cello\textquotedblright{}.) Say it once more.
\end{tcolorbox}
\section[il flauto]{il flauto — flute}

\label{lesson:IT-C516-flauto}



\begin{quote}
Before the new one: say the Italian for chairlift, then the Italian for cello.
\end{quote}

\subsection*{You'll want to know: il flauto}
\textbf{il flauto} — \textquotedblleft{}flute\textquotedblright{}.

\begin{quote}
A scuola tutti imparano il flauto dolce.
\end{quote}

\begin{quote}
\emph{At school everyone learns the recorder.}
\end{quote}

\emph{Il flauto dolce} — the recorder.

\begin{cousinweb}
\emph{flauto} comes from Old Occitan \emph{flaüt} (EtymDB); the deeper origin is not settled, though it may imitate a breath.
\end{cousinweb}

\begin{grammarlens}[title={What you've built}]
One more word for telling, explaining and describing at B1.
\end{grammarlens}

\subsection*{Your turn}
\begin{itemize}
\raggedright
  \item \emph{Say it:} \emph{il flauto}
  \item \emph{Say it:} \emph{il flauto}, once more
  \item \emph{Say it:} a sentence of your own with \emph{il flauto}
  \item \emph{Recall:} say the Italian for chairlift, then the Italian for cello, then say \emph{il flauto} again
\end{itemize}

\begin{tcolorbox}[breakable,colback=teal!4,colframe=teal!35!black,title={Before you move on}]
What does il flauto mean? (\textquotedblleft{}Flute\textquotedblright{}.) Say it once more.
\end{tcolorbox}
\section[la tromba]{la tromba — trumpet}

\label{lesson:IT-C516-tromba}



\begin{quote}
Before the new one: say the Italian for cello, then the Italian for flute.
\end{quote}

\subsection*{You'll want to know: la tromba}
\textbf{la tromba} — \textquotedblleft{}trumpet\textquotedblright{}.

\begin{quote}
Il vicino suona la tromba la domenica mattina.
\end{quote}

\begin{quote}
\emph{The neighbour plays the trumpet on Sunday mornings.}
\end{quote}

Also \emph{la tromba delle scale} — the stairwell.

\begin{cousinweb}
\emph{tromba} is usually traced to a Germanic word \emph{trumba} \textquotedblleft{}trumpet\textquotedblright{}, which may imitate the sound.
\end{cousinweb}

\begin{grammarlens}[title={What you've built}]
One more word for telling, explaining and describing at B1.
\end{grammarlens}

\subsection*{Your turn}
\begin{itemize}
\raggedright
  \item \emph{Say it:} \emph{la tromba}
  \item \emph{Say it:} \emph{la tromba}, once more
  \item \emph{Say it:} a sentence of your own with \emph{la tromba}
  \item \emph{Recall:} say the Italian for cello, then the Italian for flute, then say \emph{la tromba} again
\end{itemize}

\begin{tcolorbox}[breakable,colback=teal!4,colframe=teal!35!black,title={Before you move on}]
What does la tromba mean? (\textquotedblleft{}Trumpet\textquotedblright{}.) Say it once more.
\end{tcolorbox}
\section[la fisarmonica]{la fisarmonica — accordion}

\label{lesson:IT-C516-fisarmonica}



\begin{quote}
Before the new one: say the Italian for flute, then the Italian for trumpet.
\end{quote}

\subsection*{You'll want to know: la fisarmonica}
\textbf{la fisarmonica} — \textquotedblleft{}accordion\textquotedblright{}.

\begin{quote}
Alla festa un signore suonava la fisarmonica.
\end{quote}

\begin{quote}
\emph{At the party a man was playing the accordion.}
\end{quote}

The plural is \emph{le fisarmoniche}.

\begin{cousinweb}
\emph{fisarmonica} is a modern learned word built from Greek \emph{phŷsa} \textquotedblleft{}bellows, breath\textquotedblright{} and \emph{armonica} \textquotedblleft{}harmonica\textquotedblright{}, from Greek \emph{harmonía} \textquotedblleft{}harmony\textquotedblright{} (EtymDB).
\end{cousinweb}

\begin{grammarlens}[title={What you've built}]
One more word for telling, explaining and describing at B1.
\end{grammarlens}

\subsection*{Your turn}
\begin{itemize}
\raggedright
  \item \emph{Say it:} \emph{la fisarmonica}
  \item \emph{Say it:} \emph{la fisarmonica}, once more
  \item \emph{Say it:} a sentence of your own with \emph{la fisarmonica}
  \item \emph{Recall:} say the Italian for flute, then the Italian for trumpet, then say \emph{la fisarmonica} again
\end{itemize}

\begin{tcolorbox}[breakable,colback=teal!4,colframe=teal!35!black,title={Before you move on}]
What does la fisarmonica mean? (\textquotedblleft{}Accordion\textquotedblright{}.) Say it once more.
\end{tcolorbox}
\section[strimpellare]{strimpellare — to strum, to play badly}

\label{lesson:IT-C516-strimpellare}



\begin{quote}
Before the new one: say the Italian for trumpet, then the Italian for accordion.
\end{quote}

\subsection*{You'll want to know: strimpellare}
\textbf{strimpellare} — \textquotedblleft{}to strum, to play badly\textquotedblright{}.

\begin{quote}
Strimpello un po' la chitarra, ma non so suonare bene.
\end{quote}

\begin{quote}
\emph{I strum the guitar a bit, but I cannot play well.}
\end{quote}

Playing without much skill, for fun.

\begin{cousinweb}
The origin of \emph{strimpellare} is not settled; it is usually taken to imitate the sound of plucked strings.
\end{cousinweb}

\begin{grammarlens}[title={What you've built}]
One more word for telling, explaining and describing at B1.
\end{grammarlens}

\subsection*{Your turn}
\begin{itemize}
\raggedright
  \item \emph{Say it:} \emph{strimpellare}
  \item \emph{Say it:} \emph{strimpellare}, once more
  \item \emph{Say it:} a sentence of your own with \emph{strimpellare}
  \item \emph{Recall:} say the Italian for trumpet, then the Italian for accordion, then say \emph{strimpellare} again
\end{itemize}

\begin{tcolorbox}[breakable,colback=teal!4,colframe=teal!35!black,title={Before you move on}]
What does strimpellare mean? (\textquotedblleft{}To strum, to play badly\textquotedblright{}.) Say it once more.
\end{tcolorbox}
